\section{Petunjuk Penggunaan Buku Ajar}

\subsection{Untuk Mahasiswa}

Membaca petunjuk penggunaan buku ajar dengan saksama akan membantu mahasiswa mengarahkan pembelajaran sesuai dengan outcome yang diharapkan dan memaksimalkan manfaat dari setiap komponen OBE. Pendekatan self-directed learning yang didukung oleh buku ini mengasumsikan mahasiswa secara aktif mengikuti alur materi, mengerjakan latihan, dan melakukan refleksi terhadap pencapaian kompetensi. Petunjuk berikut dirancang untuk memandu penggunaan buku secara efektif.

Urutan materi dalam buku ini mengikuti alur logis yang direkomendasikan oleh sumber-sumber terkemuka seperti Discrete Mathematics Levin \cite{levin_discrete}. Logika proposisional dipelajari terlebih dahulu sebagai fondasi, diikuti logika predikat, metode pembuktian, induksi matematika, dan terakhir pengenalan Prolog. Banyak konsep saling bergantung; misalnya, pemahaman tentang ekuivalensi logis diperlukan sebelum menerapkan aturan inferensi, dan logika predikat memperluas notasi proposisional. Oleh karena itu, pelajarilah materi secara berurutan.

Contoh-contoh tabel kebenaran dan bukti formal disajikan untuk memperjelas konsep; disarankan mahasiswa mengerjakan contoh tersebut secara mandiri terlebih dahulu sebelum melihat solusi. Alat bantu daring seperti Natural Deduction Proof Editor dan Checker dari Open Logic Project serta Carnap dapat digunakan untuk berlatih membangun bukti dan memverifikasi kebenaran argumen \cite{proofs_checker,carnap}. Latihan di setiap bab dirancang untuk menguji penguasaan materi dan mempersiapkan mahasiswa menghadapi asesmen sumatif.

Checklist pencapaian kompetensi di akhir setiap bab berfungsi sebagai instrumen asesmen formatif dan self-assessment. Centanglah item checklist setelah yakin telah menguasainya; ini membantu mahasiswa mengidentifikasi area yang masih perlu diperdalam sebelum ujian. Berikut petunjuk penggunaan buku ajar secara ringkas:

\begin{enumerate}
  \item \textbf{Baca Sub-CPMK} di awal setiap bab untuk mengetahui kompetensi yang harus dicapai.
  \item \textbf{Pelajari materi} secara berurutan; banyak konsep saling bergantung.
  \item \textbf{Kerjakan contoh} tabel kebenaran dan bukti secara mandiri sebelum melihat solusi.
  \item \textbf{Lakukan aktivitas} pembelajaran untuk memperdalam pemahaman.
  \item \textbf{Latihan} di setiap bab untuk menguji penguasaan Anda.
  \item \textbf{Gunakan checklist} untuk self-assessment sebelum ujian.
\end{enumerate}

\subsection{Untuk Dosen}

Buku ini dapat digunakan sebagai bahan ajar utama untuk perkuliahan Logika Matematika, sumber latihan dan tugas, referensi untuk menyusun soal ujian, serta panduan untuk merancang aktivitas pembelajaran yang selaras dengan Sub-CPMK. Checklist dan komponen asesmen di setiap bab membantu dosen mengukur pencapaian CPMK mahasiswa secara terukur.
